\pdfvariable trailerid {[<C2582026083100000000000000000000><C2582026083100000000000000000000>]}
\documentclass[10pt]{article}
\usepackage[margin=0.70in]{geometry}
\usepackage{amsmath,amssymb,amsthm,booktabs,microtype,hyperref}
\hypersetup{colorlinks=true,linkcolor=blue,urlcolor=blue}
\newtheorem{theorem}{Theorem}
\newcommand{\Z}{\mathbb Z}
\title{Full-Period Mixed Congruential Dynamics:\\
A Prime-Power and CRT Classification}
\author{Route-A source-local certificate HCS-C258}
\date{31 August 2026}
\begin{document}
\maketitle
\begin{abstract}
For every $m\ge2$ we classify the affine maps
$F(x)=ax+c$ on $\Z/m\Z$ that form one cycle through all $m$ states.
The criterion is proved on prime-power quotients and assembled by the Chinese
remainder theorem (CRT).  The resulting primitive ledger and fixed counts are
exact.  This finite-ring theorem uses no target spectral data.
\end{abstract}

\section{Frozen owner and full-period theorem}
One application of $F$ is one unit of source time.  Let
\[
 L(m)=
 \begin{cases}
 \operatorname{lcm}(\operatorname{rad}m,4),&4\mid m,\\
 \operatorname{rad}m,&4\nmid m.
 \end{cases}
\]
\begin{theorem}[Hull--Dobell criterion]
The map $F$ is one $m$-cycle if and only if
\[
 (c,m)=1,\qquad p\mid a-1\quad(p\mid m),\qquad
 4\mid a-1\quad\hbox{when }4\mid m.                 \tag{1}
\]
Exactly $\varphi(m)m/L(m)$ residue pairs $(a,c)$ satisfy (1).
\end{theorem}
The classical criterion is due to Hull and Dobell~\cite{HD}; the proof below
is included to freeze the precise convention used by the certificate.

\section{Necessity and sufficiency}
If $F$ is transitive modulo $m$, every prime quotient is transitive.  When
$a\not\equiv1\pmod p$, the induced affine map has a fixed point (or is not
a permutation), which excludes a $p$-cycle.  Thus $a\equiv1\pmod p$ and
$c\not\equiv0\pmod p$.  If $4\mid m$ but $a\equiv3\pmod4$, then
$F^2(x)=a^2x+c(a+1)\equiv x\pmod4$, excluding a four-cycle.  This proves
necessity.

For sufficiency, write $p^e\Vert m$ and
\[
 S_n=1+a+\cdots+a^{n-1},\qquad
 F^n(x)-x=S_n\bigl((a-1)x+c\bigr).                   \tag{2}
\]
The second factor is a $p$-adic unit.  The standard lifting identity gives
$v_p(S_n)=v_p(n)$ for odd $p$, and also for $p=2$ under the mod-four
condition.  The $e=1$ dyadic face follows directly from oddness.  Hence every
state has exact period $p^e$ modulo $p^e$.  CRT gives exact period
$\operatorname{lcm}_{p^e\Vert m}p^e=m$, so all states comprise one cycle.
The congruences select $m/L(m)$ multipliers and $\varphi(m)$ increments.

\section{Local valuation and sharp failure faces}
The local step can be made bidirectional.  Since
$d_x=(a-1)x+c$ is a unit at every $p\mid m$, equation (2) gives
\[
 F^n(x)\equiv x\pmod {p^e}
 \quad\Longleftrightarrow\quad v_p(S_n)\ge e
 \quad\Longleftrightarrow\quad p^e\mid n.             \tag{3}
\]
For odd $p$ this is the quotient of
$v_p(a^n-1)=v_p(a-1)+v_p(n)$.  For $p=2$ and $e\ge2$,
$a\equiv1\pmod4$ gives the identical quotient; for $e=1$, the first
increment is odd and the second return gap is even.  Thus no earlier local
return is hidden by the affine shift.

The three hypotheses are individually visible.  If $p\mid(c,m)$, the
translation on the $p$-quotient is not transitive.  If
$a\not\equiv1\pmod p$, that quotient has a fixed point or loses
bijectivity.  If $4\mid m$ and $a\equiv3\pmod4$, its square is the identity
modulo four.  The theorem does not claim a full cycle decomposition after a
hypothesis fails.

\section{Primitive, zeta, and Koopman closure}
\begin{theorem}[complete source ledger]
For an admissible pair there is exactly one primitive orbit, of length $m$,
and
\[
 \#\operatorname{Fix}(F^n)=
 \begin{cases}m,&m\mid n,\\0,&m\nmid n.\end{cases}
 \qquad
 \zeta_F(t)=\exp\!\left(\sum_{n\ge1}
 \frac{\#\operatorname{Fix}(F^n)}n t^n\right)
 =\frac1{1-t^m}.                                    \tag{4}
\]
The counting-measure Koopman operator is a unitary $m$-cycle permutation.
Its eigenvalues are the $m$-th roots of unity once each, and
$\det(uI-U_F)=u^m-1$.
\end{theorem}
\begin{proof}
The full-period theorem gives one $m$-cycle.  Its $n$-th power fixes the
whole cycle exactly at multiples of $m$, proving the fixed count.  Substitution
in the exponential definition gives
$\exp(\sum_{k\ge1}t^{km}/k)=(1-t^m)^{-1}$.  A cyclic permutation is unitary
in the counting basis, and the discrete Fourier basis gives the stated
spectrum.
\end{proof}
This finite-ring affine owner differs from the repository's C172 primitive
finite-field multiplication and C204 linear rational-canonical owner:
translation, composite prime-power quotients, the mod-four exception, and
the Hull--Dobell parameter chamber are essential here.  Equation (4) remains
a source zeta, not a target Fredholm determinant.

\section{Exact certificate and boundary}
The producer and an independent implementation enumerate all $299{,}535$
affine pairs for $2\le m\le96$ with no criterion mismatch.  Six complete
cycle ledgers and odd/dyadic valuation controls are reconstructed exactly;
finite enumeration is a regression oracle rather than the all-$m$ proof.
The checker closes $300{,}210$ assertions, the symbolic implementation closes
$69$ identities, byte replay is exact, and semantic mutation rejects
$37/37$ repaired-hash changes.

The intrinsic prime-power structure supports only
\texttt{A0\_WEAK\_ARITHMETIC\_RELATION}.  There is no rational-prime
orbit dictionary, logarithmic prime clock, target divisor, or target analytic
structure.  Under \texttt{NO\_BAD\_EULER\_OR\_ROOT\_NUMBER}, the strict
tuple is
\[
\texttt{(A0\_WEAK\_ARITHMETIC\_RELATION,A1\_WEAK,A2\_FAIL,
A3\_FAIL,A4\_NATURAL\_QUANTIZATION)}.
\]
The verdict is \texttt{ROUTE\_A\_EXPLORATORY}; Route B is disabled.
The natural unitary earns only the source-local A4 coordinate: it does not
repair the failed target determinant and analytic-structure gates.

\begin{thebibliography}{9}
\bibitem{HD} T. E. Hull and A. R. Dobell,
\emph{Random number generators}, \emph{SIAM Review} \textbf{4} (1962),
230--254, \href{https://doi.org/10.1137/1004061}{doi:10.1137/1004061}.
\end{thebibliography}
\end{document}
